\chapter{Background}
\label{ch:background}
% Purpose (THESIS_GUIDE.md, Ch. 2): what a control engineer needs to read the
% thesis, organised around the three bodies of work in the goal document:
% supervisory control of grinding circuits, learned controllers, and LLM-guided
% program search. Target length 6-8 pages.

\section{Supervisory process control}
\label{sec:bg-supervisory}
Industrial processes are usually controlled hierarchically. Fast PID loops at
the base layer keep individual variables at their setpoints, while a slower
supervisory layer chooses those setpoints, coordinates interacting loops and
reacts to abnormal situations \cite{astrom2011,pomerleau2000}. In a
multivariable process, the pairing of manipulated and controlled variables
in the base layer matters: the relative gain array \cite{bristol1966} shows
whether a decentralized pairing will work with or against the process
interactions.

\section{Supervisory control of grinding circuits}
\label{sec:bg-grinding}
\TODO{The goal document's first body of work; replaces the old TODO in the
section above. Write 2-3 pages.}
% - The circuit: ball mill, sump, pump, hydrocyclone cluster, product overflow,
%   underflow returning to the mill; P80; circulating load; why it is coupled
%   and nonlinear; what is measured and what is only inferred (P80 analysers).
% - Typical regulatory loops (cyclone pressure via pump flow, sump density via
%   sump water, mill water, number of cyclones, mill speed) - keep generic here;
%   Weir's actual chain is in Chapter~\ref{ch:mill}.
% - How supervision is done in practice: operators adjusting setpoints, expert
%   systems, fuzzy-logic supervisors, MPC in mineral processing. Strengths and
%   maintenance cost of each.
% - References: e.g. Hodouin's review of process control in mineral processing,
%   Wei & Craig's survey of grinding-mill circuit control (add to references.bib).

\section{Learned controllers}
\label{sec:bg-learned}
TD-MPC \cite{hansen2022tdmpc} and TD-MPC2 \cite{hansen2024tdmpc2} learn a
latent world model, plan over a short horizon in that model and bootstrap
long-horizon value with temporal-difference learning. They are the learned,
black-box side of the comparison in this thesis.
\TODO{Expand (about 1-1.5 pages).}
% - Latent model (encoder, latent dynamics, reward and value heads), planning
%   horizon (MPPI over the latent model), value bootstrapping, reward design.
% - What it can do that a rule set cannot: an implicit model of the whole
%   circuit, trained on many episodes, plans over a horizon.
% - The reverse: inspection, guarantees on what code can do, editing by an
%   engineer; behaviour outside the training distribution is uncertain.
% - Keep Weir's specific model (autumn-night-41) for Chapter~\ref{ch:mill}.

\section{LLM-guided program search}
\label{sec:bg-program-search}
FunSearch \cite{romeraparedes2024funsearch} pairs an LLM with an automated
evaluator in an evolutionary loop: the LLM proposes programs, the evaluator
scores them, and the best programs are used as context for new proposals.
AlphaEvolve \cite{novikov2025alphaevolve} scales the same idea into a general
coding agent and reports that removing the rich problem context from the
prompt is among the most harmful ablations. Eureka \cite{ma2024eureka} uses an
LLM to write reward functions for reinforcement learning and feeds numeric
results back in a reflection step. PromptPO \cite{hatgiskessell2026promptpo}
prompts an LLM with Python descriptions of the state space, action space and
reward function and refines executable policies from rollout feedback; the
resulting policies range from tuned proportional controllers to rule-based
plans. \textcite{weng2026lbg} describes an
iterative workflow in which a policy is improved by inspecting execution
traces and rewriting code instead of through gradient updates. This thesis applies the
pattern to supervisory control, where the evaluator is a closed-loop
simulation of the process.
\TODO{One or two sentences on quality-diversity search and island models
(FunSearch's islands), since the explore phase of prompt version 6
(Section~\ref{sec:prompt-versions}) keeps three lineages for the same reason.}

\section{Specification gaming}
\label{sec:bg-gaming}
An optimizer that is scored by a metric will exploit weaknesses in that
metric \cite{amodei2016,krakovna2020}. LLM-guided search is no exception: in
this work, candidates repeatedly found ways to improve the average score
while behaving worse, which motivated the safeguards in
Section~\ref{sec:scoring}.

\section{Statistics for paired comparisons}
\label{sec:bg-stats}
\TODO{Short section (about 1 page): what the protocol in
Section~\ref{sec:protocol} relies on.}
% - Paired designs: every controller runs on the same scenarios, so compare
%   per-scenario differences.
% - Bootstrap confidence intervals; the hierarchical bootstrap (resample
%   policies, then scenarios) when the LLM side has several independent runs.
% - Significance level alpha = 0.05 (95 % intervals for H0,1 and H0,2).
% - Absence of a significant difference is not equivalence: equivalence testing
%   with two one-sided tests (TOST) and a margin fixed in advance (90 % interval
%   inside the margin). This answers Kaige's question of what "negligible" means.
% - Standardised effect size (Cohen's d_z for paired data) and why it is
%   descriptive only here.
% - Holm correction for the secondary metrics.
